Per determinare se un determinato linguaggio è regolare: mentre la dimostrazione sufficiente di AF-
decidibilità (esibendo esplicitamente un automa) può risultare noiosa e complicata, è possibile fornire dimo-
strazioni alternative (possibilmente più semplici) mostrando che l’insieme è ottenibile a partire da insiemi
AF-decidibili utilizzando operazioni “compatibili” con l’AF-decidibilità (proprietà di chiusura).
In questa sezione vedremo dunque che esistono operazioni fatte su insiemi AF-decidibili che permettono
di ottenere un altro insieme DFA-decidibile.
\begin{defi}
    Sia $A$ un alfabeto e $X_1$ e $x_2$ due inisemi di parole su $A$, ovvero $X_1,X_2 \subset A^*$. Definiamo allora:
    \begin{itemize}
    \item l'operazione di \emph{concatenazione}:
    \[
    X_1 \cdot X_2 := \{ w_1 w_2 \in A^* \mid w_i \in X_i \};
    \]

    \item l'operazione di \emph{somma}:
    \[
    X_1 + X_2 := \{ w \in A^* \mid w \in X_1 \text{ oppure } w \in X_2 \}
    \]
    (di fatto, è l'unione insiemistica di $X_1$ con $X_2$);

    \item l'operazione di \emph{star di Kleene}:
    \[
    X^* := \{ w_1 w_2 \dots w_n \mid w_i \in X, \text{for } i=1,\dots,n \text{ e } n\geq 0 \}
    \]
\end{itemize}

\end{defi}

\begin{prop}\label{prop:8}
Siano $X_1$ e $X_2$ regolari allora anche $X_1+X_2$ è regolare.


\begin{proof}
    
Sia $\mathcal{A}_1 = (T_1,q_0,F_1)$ l'automa che decide l'insieme $X_1$ e sia $\mathcal{A}_2 = (T_2,q_0',F_2)$ l'automa che decide $X_2$.
L'automa che riconosce $X_1+X_2$ è ottenuto utilizzando il prodotto cartesiano degli stati dei due automi e calcolando l'evoluzione ``parallela'' di entrambi gli automi a partire da uno stato coppia, in questo modo se una delle due componenti dello stato è accettante alla fine della computazione la parola sarà accettata poiché apparterrà ad uno dei due insiemi e quindi sarà nella loro unione.

Costruzione dell'automa $\mathcal{A} = (T,q_0'',F)$ di alfabeto $A$ e insieme degli stati $Q = Q_1\times Q_2$: definiamo
\[
q_0'' = (q_0,q_0'), \qquad F = F_1\times Q_2 \cup Q_1\times F_2
\]

e
\[
T((q_1,q_2),a) = (T_1(q_1,a), T_2(q_2,a)).
\]
Mostriamo che $\mathcal{A}$ decide precisamente $X_1+X_2$.
\[
T^*((q_0,q_0'),w)
=
\left(T_1^*(q_0,w),T_2^*(q_0',w)\right).
\]

Mostriamo ora che l'automa $\mathcal{A}$ decide l'insieme $X_1+X_2$, 
ovvero che la funzione associata all'automa coincide con la funzione caratteristica
di tale insieme.

Sia $w\in A^*$. Dalla definizione dell'insieme degli stati finali si ha:
\[
\begin{aligned}
f_{\mathcal{A}}(w)=1
&\iff
T^*((q_0,q_0'),w)\in F\\
&\iff
\left(T_1^*(q_0,w),T_2^*(q_0',w)\right)
\in (F_1\times Q_2)\cup(Q_1\times F_2)\\
&\iff
T_1^*(q_0,w)\in F_1
\ \text{oppure}\
T_2^*(q_0',w)\in F_2\\
&\iff
f_{\mathcal{A}_1}(w)=1
\ \text{oppure}\
f_{\mathcal{A}_2}(w)=1\\
&\iff
w\in X_1
\ \text{oppure}\
w\in X_2\\
&\iff
w\in X_1+X_2.
\end{aligned}
\]

Quindi, per ogni $w\in A^*$,
\[
f_{\mathcal{A}}(w)=\mathbf{1}_{X_1+X_2}(w).
\]

Pertanto l'automa $\mathcal{A}$ decide l'insieme $X_1+X_2$, e quindi
$X_1+X_2$ è regolare.
\end{proof}
\end{prop}


\begin{ex}\label{es:3}
Si consideri l'alfabeto $A=\{0,1\}$ e sia $B=\{1\}$ la restrizione all'insieme che contiene il solo simbolo $1$; si considerino l'insieme delle parole di alfabeto $A$ di lunghezza pari:
\[
X_1 = \{ w \in A^* \mid |w| = 2n,\ n\geq 0 \}
\]
e l'insieme delle parole di alfabeto $A$ che contengono un numero pari di $1$:
\[
X_2 = \{ w \in A^* \mid |w_{|B}| = 2k,\ k\geq 0 \}
\]
qui si è utilizzata la notazione $w_{|B}$ per indicare la parola $w$ ristretta ai soli elementi di alfabeto $B$, (si noti che si tratta degli insiemi considerati in Figura~\ref{fig:automa} e in Figura~\ref{fig:parita}).

Costruire l'automa che decide $X_1+X_2$.


\begin{proof}
L'automa si può ottenere utilizzando la costruzione di composizione parallela utilizzata nella prova della proposizione a partire dai due automi per $X_1$ e $X_2$.
\end{proof}
\end{ex}


Mostriamo ora un'altra proprietà di chiusura degli insiemi regolari: la classe degli insiemi regolari è chiusa per concatenazione, ovvero:

\begin{prop}\label{prop:9}
Siano $X_1$ e $X_2$ regolari allora anche $X_1\cdot X_2$ è regolare.
\end{prop}

Potremmo tentare la costruzione di un automa come nella proposizione precedente in modo che ogni parola $w\in X_1X_2$ abbia un inizio riconosciuto dall'automa che riconosce $X_1$ e continui con una seconda parte che sia accettata dall'automa che riconosce $X_2$. In pratica si potrebbe concatenare il primo automa con il secondo, ma il problema sarebbe che potrebbero esserci più modi di fattorizzare $w$ in modo che il primo automa accetti la prima parte di $w$. Per poter definire l'automa concatenazione bisognerebbe sapere a quale punto passare dal primo al secondo automa (e non è detto che ci sia un unico punto che vada bene per tutte le parole di $X_1\cdot X_2$). Per risolvere questo problema si introduce una nuova nozione: la nozione di computazione \emph{non-deterministica}.